\chapter*{How This Volume Is Organised}
\addcontentsline{toc}{chapter}{How This Volume Is Organised}
\label{ch:plan}

\section*{What this volume solves}
\addcontentsline{toc}{section}{What this volume solves}

\emph{\term{Arrangement and Orchestration}{arrangement}} answers one practical question: \textbf{given a melody, a chord progression, or a drum groove, how do you turn it into a complete, layered piece that can be realised in Yeban?} It is about what to write, for which voice, and on which layer. It is not about polishing recorded sound; that belongs to the mixing and mastering volume.

The volume assumes you can create a project, add tracks, place clips, and open the piano roll in Yeban. If not, read the User Guide chapters on projects, tracks and clips, and the piano roll first.

\section*{Map of the eleven chapters}
\addcontentsline{toc}{section}{Map of the eleven chapters}

\begin{longtable}{p{0.07\textwidth}p{0.28\textwidth}p{0.32\textwidth}p{0.25\textwidth}}
\toprule
Ch. & Topic & After studying & Where it lands in Yeban \\
\midrule
\endfirsthead
\toprule
Ch. & Topic & After studying & Where it lands in Yeban \\
\midrule
\endhead
1 & Arrangement overview & Use a four-layer role card to hear what a song lacks & Six-track demo, solo/mute, piano roll \\
2 & Song structure & Draw the energy curve of verse, chorus and bridge & Section markers, arrangement area, overview panel \\
3 & Rhythm section & Make drums, bass and comping share the groove & Audio and MIDI tracks, velocity, quantise \\
4 & Harmony layer & Write pads, arpeggios and chordal instruments without mush & Chord lane, piano roll, device rack \\
5 & Melody and counterpoint & Write a lead, a counter-line and answering phrases & MIDI clip copies, transpose, velocity lane \\
6 & String orchestration & Divide a string section into sensible voices & String tracks, SFZ instruments, automation \\
7 & Woodwind and brass & Choose registers and combinations without covering the vocal & Multi-track layering, pan, faders \\
8 & Electronic and pop arranging & Build texture with synth layers, samples and loops & Synth parameters, asset pool, clip loops \\
9 & Vocal arranging & Arrange lead, harmony and stacked vocals & Multiple vocal tracks, clip alignment, stem export \\
10 & Walkthroughs & Carry a folk and a pop example from start to finish & Project save, WAV export, stem export \\
11 & Genre and orchestration & Build eight real-world genre skeletons and judge instrumentation with four rulers & Track/section markers, piano roll, automation, Solo/Mute \\
\bottomrule
\end{longtable}

\section*{Three learning paths}
\addcontentsline{toc}{section}{Three learning paths}

\begin{itemize}
\item \textbf{From zero}: read Chapters 1--5 in order and do at least one eight-bar exercise per chapter; finish with the walkthrough in Chapter 10, then pick one unfamiliar genre in Chapter 11 and build its eight-bar skeleton.
\item \textbf{Producer}: start with Chapters 2, 3, 4, 8, 10 and 11; consult Chapters 6 and 7 when a string or brass part appears.
\item \textbf{Orchestrator}: read Chapters 1, 2, 4, 6, 7, 9 and 11 closely; use Chapters 3 and 5 as support for rhythm and melody.
\end{itemize}

\section*{Writing and review conventions}
\addcontentsline{toc}{section}{Writing and review conventions}

\begin{itemize}
\item A term is introduced as ``Chinese (English)'' on first use and indexed at the end. New terms for this volume are listed in the term table below.
\item UI screenshots use the placeholder \verb|\screenshot{name}{caption}|. Names are lowercase and hyphenated; management captures and fills them from the receipt list.
\item Concept figures, flow diagrams and notation examples are drawn in TikZ and shipped with the source. Text in a figure explains relationships; it never pretends to be a real screenshot.
\item Each chapter ends with a help-topic mapping table from section to in-app help anchor. Anchors are lowercase and hyphenated for embedding into the application help.
\item Product claims are limited to entries that exist on the live tip and in the current in-app manual. Unimplemented or planned features are omitted or explicitly marked ``not yet available''.
\end{itemize}

\section*{Glossary of this volume}
\addcontentsline{toc}{section}{Glossary of this volume}

This table lists the specialist terms used in this volume and their meaning here. The entries have been merged into the shared glossary (\texttt{common/glossary-en.tex}); duplicates and conflicts skipped during the merge are recorded in \texttt{GLOSSARY-MERGE.md}.

\begin{longtable}{p{0.24\textwidth}p{0.24\textwidth}p{0.44\textwidth}}
\toprule
Chinese & English & Meaning in this volume \\
\midrule
\endfirsthead
\toprule
Chinese & English & Meaning in this volume \\
\midrule
\endhead
编曲层次 & arrangement layer & One of four functional layers: rhythm, harmony, melody, colour/atmosphere; not a fixed track count. \\
角色卡 & role card & Three questions answered before writing a part: function, register, rhythmic character. \\
能量曲线 & energy curve & A bar-by-bar graph of loudness, density and instrument count across sections. \\
织体 & texture & How voices combine in register, rhythmic density and timbre. \\
声部避让 & voice avoidance & Keeping two parts from competing in the same register and on the same beat. \\
节奏组 & rhythm section & Drums, bass and comping instruments responsible for groove and root motion. \\
铺底 & pad bed & Sustained or slow-moving chords that carry the harmonic background. \\
琶音 & arpeggio & A chord played as a sequence of single notes rather than a block. \\
对位 & counterpoint & Two or more melodic lines that move independently yet fit together. \\
呼应句 & call-and-response & A short answering phrase after a lead phrase, usually in another voice. \\
弦乐组 & string section & First and second violins, viola, cello and double bass as a divided group. \\
木管组 & woodwind section & Flute, oboe, clarinet, bassoon and related woodwinds as a group. \\
铜管组 & brass section & Trumpet, horn, trombone, tuba and related brass as a group. \\
叠唱 & vocal stacking & Recording the same lyric several times to make a thicker vocal sound. \\
走查 & walkthrough & A section-by-section review of a complete piece against the Yeban project. \\
Verse & verse & The section that tells the story and gives information; usually lower energy. \\
Chorus & chorus & The most memorable section, with the fullest harmony and instrumentation. \\
Pre-chorus & pre-chorus & The lift between verse and chorus; its job is to create expectation. \\
Bridge & bridge & A contrasting middle section, often changing chords, groove or timbre; not a second verse. \\
Outro & outro & The closing or fading section. \\
Section marker & section marker & A mark on the arrangement ruler that shows a section's start and end. \\
Instrumentation budget & instrumentation budget & A per-section table of what each layer adds or removes, controlling energy and contrast. \\
Transition & transition & A device joining two sections: hard cut, fill, harmonic transition or space. \\
Hard cut & hard cut & Changing section directly on a bar line or strong beat without a fill. \\
Fill & fill & A one-bar drum, bass or guitar gesture that carries the ear into the next section. \\
Groove & groove & The rhythmic character created by accents, space and drive. \\
Anchor & anchor & The voice or placement that pins the beat and root, usually kick and bass roots. \\
Backbeat & backbeat & The snare accent on 2 and 4 or other weak beats. \\
Half-time & half-time & A drum feel that stretches the perceived time by roughly a factor of two. \\
Four-on-the-floor & four-on-the-floor & A dance skeleton with a kick on every beat. \\
Root & root & The bass note of a chord; the bass must make it audible first. \\
Same-note repetition rate & same-note repetition rate & The share of notes in a bass or groove part that repeat the same pitch. \\
Leap & leap & A large interval jump between adjacent notes; this volume counts five semitones or more. \\
Comping & comping & Rhythmic guitar/keyboard writing that gives chords on key beats and leaves space elsewhere. \\
Quantize grid & quantize grid & The timing grid notes align to; align the skeleton, humanise the ornaments. \\
Harmony layer & harmony layer & The functional layer carrying chord motion, texture and harmonic support. \\
Harmonic rhythm & harmonic rhythm & The rate of chord change, e.g. once per bar or every two beats. \\
Voice leading & voice leading & How individual voices move between adjacent chords; smoothness first. \\
Close voicing & close voicing & A chord with its tones packed close, suitable for mid/high registers. \\
Open voicing & open voicing & A chord with its tones spread out, suitable for pads and strings. \\
Inversion & inversion & Moving the root away from the lowest voice, using the third, fifth or seventh as bass. \\
Top voice & top voice & The highest voice of a chord voicing, often used as a melodic line. \\
Chord lane & chord lane & The lane below the arrangement area that shows and aligns chord changes. \\
Tension & tension & Unresolved instability in a chord, often from sevenths/ninths or delayed resolution. \\
Pedal & pedal & Holding one note across chord changes to create connection or suspense. \\
Motif & motif & The shortest recognisable melodic unit, usually 2--5 notes. \\
Phrase & phrase & A melodic unit with a beginning and a landing, commonly 4 or 8 bars. \\
Antecedent & antecedent & The first half of a phrase pair that poses the question and builds expectation. \\
Consequent & consequent & The second half that answers the antecedent and resolves expectation. \\
Contour & contour & The melodic shape in pitch: rising, falling, arch, etc. \\
Target note & target note & The chord tone a phrase lands on, deciding stability, suspension or resolution. \\
Counter-melody & counter-melody & An independent melodic line running with the lead in a different register or rhythm. \\
Variation & variation & Developing a motif by transposition, rhythmic change or direction while keeping its identity. \\
Breath & breath & The pause between phrases that makes the sentence boundary audible. \\
Stress & stress & The emphasised syllable or beat in a lyric or melody. \\
String orchestration & string orchestration & Writing the string section as five parts with clear roles and registers. \\
Divisi & divisi & Splitting one string section into two halves to enrich harmony. \\
Bowing & bowing & Legato, pizzicato, short notes and other bowing styles that decide phrase character. \\
Legato & legato & Notes connected without a break; suited to pads and long lines. \\
Pizzicato & pizzicato & Short, granular string writing suited to light passages. \\
Tutti & tutti & The density of all five string parts sounding together. \\
Chamber-like & chamber-like & Sparse string writing with three parts or fewer. \\
Bass anchor & bass anchor & The root and strong-beat landing held by cello/double bass. \\
Voice crossing & voice crossing & A lower part moving above a higher one, blurring the roles. \\
Voicing & voicing & The distribution and spacing of chord tones across the parts. \\
Woodwind orchestration & woodwind orchestration & Writing flute, oboe, clarinet and bassoon as separate lines. \\
Brass orchestration & brass orchestration & Layering trumpet, horn, trombone and tuba with accents. \\
Timbral colour & timbral colour & The brightness, warmth and emotional tendency a timbre brings. \\
Projection & projection & How strongly an instrument cuts through at equal volume; not the same as fader size. \\
Brass stab & brass stab & A short brass chord at a section change or the chorus downbeat. \\
Family grouping & family grouping & Pairing woodwinds with woodwinds, or brass with brass. \\
Mixed grouping & mixed grouping & Using woodwinds and brass together for contrast. \\
Doubling & doubling & Having another instrument repeat a lead/part, usually in a new register or interval. \\
Woodwind duet & woodwind duet & Two woodwinds each playing an independent line. \\
Brass choir & brass choir & A brass block arranged high--middle--middle-low--low. \\
Electronic timbre layer & electronic timbre layer & Splitting an electronic sound into sub, body, air and rhythm functions. \\
Sub layer & sub layer & The synth layer carrying the lowest band and weight. \\
Body layer & body layer & The midrange synth layer that carries the timbre's identity. \\
Air layer & air layer & The layer handling high detail, sheen and phrase-end decoration. \\
Timbre layering & timbre layering & Giving each layer an independent envelope and time shape. \\
Frequency-band division & frequency-band division & Assigning low, body, air and rhythm to separate bands so they do not mask each other. \\
Envelope & envelope & Attack/Decay/Release and related parameters that shape the note's start and tail. \\
Filter & filter & A device using cutoff/resonance to control brightness and peaks. \\
Asset pool & asset pool & Where imported audio material is stored before being placed on an audio track. \\
Sample & sample & Existing recorded material imported and placed on an audio track. \\
Vocal arrangement & vocal arrangement & Assigning lead, double, harmony, stack and ad-lib to roles. \\
Lead vocal & lead vocal & The vocal layer singing the main lyric, carrying narrative and emotion. \\
Double & double & Singing the same lyric again and layering it on the lead. \\
Vocal harmony & vocal harmony & A second melodic vocal line a third/sixth below the lead. \\
Ad-lib & ad-lib & A short response, shout or breath at the end of a lead phrase. \\
Vocal layer & vocal layer & One functional layer within a vocal arrangement. \\
Stack layer count & stack layer count & The number of stacked layers of the same lyric, usually three to five. \\
Lyric intelligibility & lyric intelligibility & Whether every word is clear; the first criterion of vocal arrangement. \\
Lead-centered & lead-centered & Placing the lead at the centre and spreading other layers to the sides. \\
Stem export & stem export & Exporting each track as a separate WAV. \\
Short song & short song & A teaching demo around sixteen bars, not a commercial length. \\
Demo project & demo project & A teaching project built by this book's methods, not a preinstalled repository project. \\
Walkthrough checklist & walkthrough checklist & A section-by-section check of entrances, energy, low end, lead, colour and delivery. \\
Track chart & track chart & A table listing track, type and assigned role. \\
Section chart & section chart & A table listing bars, sections and arrangement actions. \\
Export receipt & export receipt & The report after WAV/stem export with frames, true peak, counts and skipped items. \\
Atomic write & atomic write & Saving that never leaves half a project even if power fails midway. \\
Version time machine & version time machine & A panel that draws the project commit chain and returns to any version. \\
Bus limiter & bus limiter & A WAV export switch that peak-limits the master bus to prevent clipping. \\
Diagnostic report bundle & diagnostic report bundle & A troubleshooting package, not an audio deliverable. \\
曲风 & genre & A set of instrumentation, groove, harmony and form choices, not a timbre label. \\
编制 & instrumentation & Which instruments/voices a piece uses and how they divide the work. \\
律动型 & groove pattern & The repeating time skeleton of drums, bass and comping. \\
循环动机 & riff & A short repeated instrumental figure that becomes a memory point. \\
钩子 & hook & The most memorable short phrase, timbre or rhythm. \\
主调织体 & homophony & One lead melody with chordal accompaniment; clear hierarchy. \\
复调织体 & polyphony & Two or more independent melodies moving at once. \\
织体分层 & stratified texture & Low, body, air and rhythm each occupy their own layer. \\
声部密度 & voice density & How many independent parts are saying different things at once. \\
音域 & range & The lowest to highest note an instrument or voice can produce. \\
移调 & transposition & Converting written to sounding pitch, or moving a whole line. \\
适奏音区 & tessitura & The most comfortable, most-used part of an instrument's or voice's range. \\
留白 & space & A deliberate gap that makes density and contrast audible. \\
固定音型 & ostinato & A short figure repeated, varied by displacement, ornament or timbre. \\
升压段 & build-up & The lift before an EDM drop. \\
电子落点 & drop & The section where EDM low end and body return together. \\
侧链 & sidechain & Kick-triggered compression that briefly ducks pad/bass low end. \\
摇摆 & swing & Long--short eighth-note feel created by delaying the second half. \\
行走贝斯 & walking bass & One note per beat, connecting roots with passing tones. \\
强和弦 & power chord & Root plus fifth, common in rock guitar. \\
蓝调音 & blue note & Bent lowered third/seventh with blues colour. \\
影视配乐 & film score & Instrumental music organised by picture, narrative, motives and orchestration. \\
国风 & guofeng & A style marked by pentatonic scales, Chinese timbres, ornament and breath. \\
情绪曲线 & emotional curve & A section-by-section map of emotional intensity and release points. \\
耳糖 & ear candy & A small timbre or phrase that carries no main narrative but makes a section taste better. \\
流行 & pop & A song type whose reach centres on hook, groove, contrast and payoff. \\
托普线 & topline & The main melody at the top of a song, usually carried by the lead vocal. \\
失真吉他 & distorted guitar & A guitar timbre that uses distortion/overdrive to thicken harmonics and add mid-low bite. \\
静响对比 & quiet-loud dynamics & Section contrast created by density, distortion and force: hold back, then release. \\
吉他墙 & guitar wall & Several distorted guitars stacked into a wide, thick chord backdrop, often in the chorus. \\
摇滚 & rock & A song type driven by riff, backbeat, distortion and release. \\
电子舞曲 & electronic dance music & A dance genre built on four-on-the-floor, frequency-band division and the anticipation-release arc of the drop. \\
泵感 & pumping & The breathing feel created when other layers duck on each kick and return afterwards. \\
截止频率 & cutoff frequency & The point above which a filter attenuates; raising it makes a sound brighter. \\
重复中的微变 & variation within repetition & Writing a short motif that returns with only one small detail changed. \\
\bottomrule
\end{longtable}

\section*{What this volume does not promise}
\addcontentsline{toc}{section}{What this volume does not promise}

This volume does not replace mixing, mastering or plug-in manuals. It does not claim that one instrumentation works for every song, and it does not present ``professional arranging'' as a single correct recipe. Every register, frequency and velocity number is a \textbf{starting point}, not a law. Use it to remove obvious mistakes, then judge with your ears and reference records.

\begin{helpmapping}
\begin{longtable}{p{0.45\textwidth}p{0.45\textwidth}}
\toprule
Section & Help anchor \\
\midrule
How this volume is organised & \texttt{arrangement-plan} \\
Map of the eleven chapters & \texttt{arrangement-chapter-map} \\
Three learning paths & \texttt{arrangement-learning-paths} \\
Writing and review conventions & \texttt{arrangement-conventions} \\
New terms in this volume & \texttt{arrangement-glossary} \\
\bottomrule
\end{longtable}
\end{helpmapping}
